\chapter{Faraday's law, inductance, RL and RLC}\label{ch:faraday}

The last symmetry to break: changing magnetic flux \emph{creates}
electric field --- a non-conservative one whose loops replace the
conservative fields of chapters 9--11. Motors, generators, transformers
and every wireless charger are this chapter monetised.

\section{Flux and motional EMF: the sliding rod first}\label{sec:far-rod}

A rod sliding on rails through $\mathbf B$ (figure~\ref{fig:rod}) sweeps
area: $\Phi_B=BLx$, so $\mathcal E=-d\Phi_B/dt=-BLv$. The same answer
comes from Lorentz force on the rod's charges --- motional EMF is
Faraday's law's first instance.

\begin{figure}[htbp]\centering
\begin{tikzpicture}[scale=1.0,font=\scriptsize]
\draw[very thick] (-0.3,0) -- (3.3,0);
\draw[very thick] (-0.3,1.2) -- (3.3,1.2);
\draw[very thick,accent] (1.6,-0.1) -- (1.6,1.3);
\draw[->,thick,amber] (1.6,0.6) -- (2.5,0.6) node[right]{$v$};
\foreach \x in {0.4,0.9,2.4,2.9}{\foreach \y in {0.4,0.9}{\node at (\x,\y) {$\times$};}}
\node[align=left] at (1.6,-0.7) {$\mathcal E=BLv$; Lenz pushes back\\$F=ILB$ against $v$};
\end{tikzpicture}

\caption{Sliding rod: the swept flux gives $\mathcal E=BLv$, and Lenz's
law retards the motion.}
\label{fig:rod}
\end{figure}

\begin{example}[Sliding rod power audit]\label{ex:rod}
$B=0.8$\,T, $L=0.5$\,m, $R=4\,\Omega$, pulled at $v=3$\,m/s. Find
$\mathcal E$, $I$, the retarding force, and audit the powers.
\par\smallskip
\textbf{Strategy.} $\mathcal E=BLv$; $F=ILB$ opposes $v$ by Lenz.
\begin{solution}
$\mathcal E=1.20$\,V, $I=0.30$\,A, $F=0.12$\,N. Mechanical power in
$Fv=0.360$\,W; electrical generated $\mathcal EI=0.360$\,W; resistor
burns $I^{2}R=0.360$\,W. Three ways, one number --- the code prints
all of them. \checkmark
\end{solution}
\end{example}

\section{Faraday's law}\label{sec:far-law}

\begin{equation}
  \mathcal E=-\frac{d\Phi_B}{dt}
  =-\frac{d}{dt}\int\mathbf B\cdot d\mathbf A,\label{eq:far}
\end{equation}
counting every way flux can change: move, rotate, grow the loop or grow
$\mathbf B$ (figure~\ref{fig:flux4}). The minus sign is Lenz ---
the next section's whole job.

\begin{figure}[htbp]\centering
\begin{tikzpicture}[scale=0.55,font=\scriptsize]
\draw[very thick] (0,0) rectangle (2.4,1.6);
\foreach \x in {0.6,1.2,1.8}{\node at (\x,0.8) {$\times$};}
\draw[->,thick,amber] (2.4,0.8) -- (3.1,0.8);
\node[align=left] at (1.2,-0.6) {move};
\draw[very thick,rotate around={20:(5.4,0.8)}] (4.4,0) rectangle (6.4,1.6);
\foreach \x/\y in {4.9/0.7,5.5/0.9,6.0/1.1}{\node at (\x,\y) {$\times$};}
\node[align=left] at (5.4,-0.6) {rotate};
\draw[very thick] (7.8,0) circle (0.5);
\draw[very thick,dashed] (7.8,0) circle (0.95);
\node at (7.8,0.5) {$\times$};
\node[align=left] at (7.8,-1.3) {grow};
\draw[very thick] (10.4,0) rectangle (12.0,1.6);
\foreach \x/\y in {10.9/0.8,11.4/0.8}{\node at (\x,\y) {$\times$};}
\node at (11.6,0.8) {$\times$};
\node[align=left] at (11.2,-0.6) {grow B};
\end{tikzpicture}

\caption{Four ways to change flux: move, rotate, grow the loop, or
change $\mathbf B$ itself.}
\label{fig:flux4}
\end{figure}

\section{Lenz's law: direction by conservation}\label{sec:far-lenz}

The induced current's field opposes the \emph{change} --- not the
field. If it aided the change you'd get a runaway free-energy machine;
Lenz's minus is energy conservation wearing a sign. Applications:
the rod's retarding force, eddy braking (figure~\ref{fig:eddy}),
and every inductor's stubbornness.

\begin{figure}[htbp]\centering
\begin{tikzpicture}[scale=0.9,font=\scriptsize]
\draw[very thick] (-2.2,-1) rectangle (2.2,1);
\foreach \x in {-1.4,-0.7,0,0.7,1.4}{\node at (\x,0) {$\times$};}
\draw[->,thick,amber] (0,-1.4) -- (0,-0.7);
\node[align=left] at (0,-1.8) {magnet falls into tube:\\eddies oppose the fall};
\draw[accent,dashed,->] (-0.5,0.3) arc (30:330:0.45);
\draw[accent,dashed,->] (0.6,-0.1) arc (30:330:0.45);
\end{tikzpicture}

\caption{Eddy braking: a falling magnet's induced currents fight its
descent --- magnetic drag.}
\label{fig:eddy}
\end{figure}

\begin{example}[Lenz on entry and exit]\label{ex:lenz}
A square loop enters a uniform-$B$ region at constant $v$, crosses,
and exits. Describe the current and the required pulling force at each
stage.
\par\smallskip
\textbf{Strategy.} Only \emph{changes} count: entry, transit, exit.
\begin{solution}
Entry: flux grows $\Rightarrow$ opposing field $\Rightarrow$ current
counterclockwise; force opposes entry, magnitude
$B^{2}L^{2}v/R$. Transit: $\Phi$ constant, no current, no force.
Exit: flux falls $\Rightarrow$ field-adding current $\Rightarrow$
clockwise; same retarding magnitude. The loop is dragged \emph{in both
directions} --- braking on the way in and on the way out. \checkmark
\end{solution}
\end{example}

\section{Induced $\mathbf E$: the non-conservative leap}\label{sec:far-e}

Faraday's law written locally says
$\oint\mathbf E\cdot d\boldsymbol\ell=-d\Phi_B/dt$: an electric field
that circulates, whose ``potential'' can't exist (the loop integral is
nonzero). Around a changing solenoid the field curls in rings
(figure~\ref{fig:sole}) --- the betatron principle, and the conceptual
break from every $\mathbf E$ in chapters 9--11.

\begin{figure}[htbp]\centering
\begin{tikzpicture}[scale=0.9,font=\scriptsize]
\foreach \x in {-1.4,-0.7,0,0.7,1.4}{\draw[very thick] (\x,0.4) ellipse (0.1 and 0.22);
\draw[very thick] (\x,-0.4) ellipse (0.1 and 0.22);}
\foreach \x in {-1.4,-0.7,0,0.7}{\draw[very thick] (\x,0.4) -- (\x+0.7,0.4);
\draw[very thick] (\x,-0.4) -- (\x+0.7,-0.4);}
\node at (0,0) {$B$ increasing $\to$};
\foreach \r in {0.9,1.3}{\draw[accent,dashed,->] (0:\r) arc (0:270:\r);}
\node[align=left] at (0,-1.8) {induced $\mathbf E$ circulates\\around the changing flux};
\end{tikzpicture}

\caption{Induced $\mathbf E$ circulates around changing flux: no
potential describes it.}
\label{fig:sole}
\end{figure}

\section{Inductance: a circuit element's inertia}\label{sec:far-ind}

A solenoid's own flux links itself:
$L=\mu_0n^{2}A\ell$ (from Amp\`ere's $B=\mu_0nI$ and
$\Phi$ per turn); its EMF is $-L\,dI/dt$. Mutual inductance is the same
number between two circuits --- the transformer's measure.

\begin{example}[Solenoid inductance end-to-end]\label{ex:solL}
Derive $L$ for a solenoid of $N$ turns, length $\ell$, area $A$.
\par\smallskip
\textbf{Strategy.} $B$ by Amp\`ere, flux linkage, divide by $I$.
\begin{solution}
$B=\mu_0nI$ gives flux per turn $\Phi=BA=\mu_0nIA$; linkage over all
turns: $N\Phi=\mu_0n^{2}A\ell I$, so
$L=N\Phi/I=\mu_0n^{2}A\ell$. Numbers: $n=1000$/m, $A=2$\,cm$^{2}$,
$\ell=10$\,cm give $L=0.25$\,mH. \checkmark
\end{solution}
\end{example}

\begin{example}[RL energisation]\label{ex:rl}
A $50$\,mH inductor energises through $10\,\Omega$ from $12$\,V.
Find $\tau$, the initial slope $dI/dt$, and $I_\infty$.
\par\smallskip
\textbf{Strategy.} $\mathcal E-L\,dI/dt=IR$; same ODE family as RC.
\begin{solution}
$\tau=L/R=5$\,ms; at $t=0$ the inductor acts open so
$dI/dt=\mathcal E/L=240$\,A/s; $I_\infty=1.2$\,A.
Current climbs the same $1-e^{-t/\tau}$ curve that charge did in
chapter~12 (figure~\ref{fig:rlcur}). \checkmark
\end{solution}
\end{example}

\begin{figure}[htbp]\centering
\begin{tikzpicture}
\begin{axis}[width=8.2cm,height=4.5cm,xlabel={$t/\tau$},ylabel={},
xmin=0,xmax=6,ymin=0,ymax=1.2,font=\small,legend pos=south east]
\addplot[very thick,accent,domain=0:6,samples=60] {1-exp(-x)};
\addlegendentry{$i/I_\infty$}
\addplot[very thick,amber,domain=0:6,samples=60] {exp(-x)};
\addlegendentry{$V_L/\mathcal E$}
\addplot[dashed] coordinates {(1,0) (1,0.632)};
\end{axis}
\end{tikzpicture}

\caption{RL energisation: $i$ approaches $I_\infty$ on $\tau=L/R$ while
$V_L$ decays --- the mirror of RC.}
\label{fig:rlcur}
\end{figure}

\section{Stored energy and the field density}\label{sec:far-u}

Charging an inductor to $I$ stores $U=\tfrac12LI^{2}$; per volume the
magnetic energy density is $u_B=B^{2}/2\mu_0$ --- twin to the electric
$u_E=\varepsilon_0E^{2}/2$ of chapter~11.

\section{LC and RLC: the electric oscillator}\label{sec:far-lc}

With no battery the charge sloshes: $U_E\leftrightarrow U_B$ at
$\omega_0=1/\sqrt{LC}$ (figure~\ref{fig:slosh}) --- equation
\eqref{eq:shm} of chapter~7 with $m\to L$, $k\to1/C$, $b\to R$.
Damped RLC rings at $\omega_d=\sqrt{\omega_0^{2}-\gamma^{2}/4}$ with
$\gamma=R/L$: every result of \S\ref{sec:osc-driven} carries over.

\begin{figure}[htbp]\centering
\begin{tikzpicture}[scale=0.85,font=\scriptsize]
\draw[very thick] (-2.6,-1.2) rectangle (-0.6,1.2);
\draw[thick] (-1.9,1.2) -- (-1.9,1.6);\draw[thick] (-1.7,1.2) -- (-1.7,1.4);
\node at (-1.6,0) {$U_E$};
\node[align=left] at (-1.6,-1.7) {capacitor:\\charged};
\draw[very thick] (0.6,-1.2) rectangle (2.6,1.2);
\draw[thick,decorate,decoration={coil,amplitude=2pt,segment length=6pt}] (1.3,-1.2) -- (1.3,1.2);
\node at (1.6,0) {$U_B$};
\node[align=left] at (1.6,-1.7) {inductor:\\current max};
\draw[->,very thick,accent] (-0.6,0.2) -- (0.6,0.2);
\draw[->,very thick,accent] (0.6,-0.2) -- (-0.6,-0.2);
\node[align=left] at (0,-2.4) {$U_E+U_B$ sloshes at $\omega_0=1/\sqrt{LC}$};
\end{tikzpicture}

\caption{The LC tank: energy trades capacitor $\leftrightarrow$
inductor, quarter period apart.}
\label{fig:slosh}
\end{figure}

\begin{example}[RLC ringing vs prediction]\label{ex:rlc}
$L=50$\,mH, $C=100$\,\textmu F, $R=2\,\Omega$: predict $\omega_0$,
$\gamma$, $\omega_d$ and the decay per cycle.
\par\smallskip
\textbf{Strategy.} The oscillator formulas verbatim.
\begin{solution}
$\omega_0=1/\sqrt{LC}=447$\,rad/s, $\gamma=R/L=40$\,s$^{-1}$,
$\omega_d=\sqrt{\omega_0^{2}-\gamma^{2}/4}=446.8$\,rad/s;
decay per cycle $e^{-\gamma T/2}=0.755$ --- the code's measured trace
matches both numbers exactly. \checkmark
\end{solution}
\end{example}

\begin{lstlisting}[language=Python,caption={\texttt{code/ch14\_rlc.py} --- rod audit + RLC ringing vs theory.},label={lst:rlc}]
d2q = -(R/L)*dq - q/(L*C)
\end{lstlisting}

Output of a real run:
\begin{lstlisting}[numbers=none]
rod: emf=BLv=1.200 V, I=0.300 A, retarding F=0.120 N
  mech power F*v=0.360 W  |  elec power emf*I=0.360 W  |  I^2 R=0.360 W

RLC: w0=447.2, gamma=40.0, wd_theory=446.8
maxima: [('0.0141s', '7.55e-04'), ('0.0281s', '5.70e-04'), ('0.0422s', '4.30e-04')]
wd_measured = 446.8 (theory 446.8)
decay ratio/cycle = 0.755 (theory e^(-gT/2)=0.755)
\end{lstlisting}

The audit closes to the watt and the ring frequency/decay match theory
to all printed digits --- chapter~7's oscillator, here in copper and
charge.

\section{In engineering practice: the power grid's physics}\label{sec:far-eng}

Generators are figure~\ref{fig:flux4}'s rotation panel run backwards
(work in, EMF out); transformers are mutual inductance with an iron
guide; eddy braking stops trains and roller-coasters without contact;
and induction cooktops are \S\ref{sec:far-e} making currents in your
pot --- the induced-$\mathbf E$ field turned appliance.

\section{Common mistakes}\label{sec:far-mistakes}

\begin{tripwire}[EMF sign vs magnitude conflated]
Faraday gives magnitude; Lenz gives direction. Solve them separately
then combine --- deciding both at once is where sign errors breed.
\end{tripwire}
\begin{tripwire}[Induced $\mathbf E$ treated as conservative]
It has no potential: $\oint\mathbf E\cdot d\boldsymbol\ell\ne0$. All
the potential bookkeeping of chapter~11 fails here.
\end{tripwire}
\begin{tripwire}[$\Phi_B$ vs $B$ confusion]
$\Phi$ integrates area; EMF needs its \emph{rate}. A constant huge
field induces nothing.
\end{tripwire}
\begin{tripwire}[Motional + transformer EMF double-counted]
They are two descriptions of the same physics --- count one.
\end{tripwire}
\begin{tripwire}[``Inductor resists current'' myth]
It resists \emph{changes}: at steady state it's a wire; only the
transient feels its inertia.
\end{tripwire}

\section{Summary}\label{sec:far-summary}

\begin{table}[ht]\centering
\begin{tabularx}{\linewidth}{@{}lY@{}}
\toprule
\tablehead\thead{Idea} & \thead{One-line version}\\
\midrule
Faraday & $\mathcal E=-d\Phi_B/dt$; four ways to change flux \\
Lenz & induced current opposes the change --- conservation's sign \\
induced E & circulates; no potential \\
inductor & $L=\mu_0n^{2}A\ell$; $\tau=L/R$; $U=\tfrac12LI^{2}$ \\
LC/RLC & $\omega_0=1/\sqrt{LC}$; ch7 oscillator in disguise \\
\bottomrule
\end{tabularx}
\caption{Induction cheat sheet.}
\label{tab:farsummary}
\end{table}

\begin{wins}[Definition of done]
You can now: compute motional and transformer EMF, direct Lenz by
conservation, derive $L$ of a solenoid, solve RL/LC transients, and run
the rod power audit and RLC ringing checks.
\end{wins}

\chapternotes{%
The motional-first order follows MIT 8.02SC modules 21--25
\cite{ocw801} and AP Physics C E\&M unit 13 \cite{apced}; Griffiths
ch.~7 is the deeper anchor (in print); PhET's Faraday app the visual
one. The numbers are the listing's real output.}

\begin{rulebox}[{Watch instead}]
\begin{enumerate}
\item \emph{Electromagnetic Induction, Faraday's Law, Lenz Law} ---
Walter Lewin, MIT 8.02x Lecture 16, 51:24,
\url{https://youtu.be/nGQbA2jwkWI} ---
\S\ref{sec:far-rod}--\S\ref{sec:far-e} at full length (accessed
2026-10-11).
\item \emph{Faraday's Law Introduction} --- Khan Academy, 5:36,
\url{https://youtu.be/vcStzn55MG0} --- \S\ref{sec:far-law} worked
quickly (accessed 2026-10-11).
\end{enumerate}
\end{rulebox}

\section*{Further reading}
\addcontentsline{toc}{section}{Further reading}

\begin{itemize}
\item MIT OCW \emph{8.02SC}, induction modules 21--25. At
\url{https://ocw.mit.edu}, course ``8.02SC'' (verified 2026-10-11).
\item OpenStax \emph{College Physics 2e}, ch.~23 ``Electromagnetic
Induction'' --- free text: \url{https://openstax.org} (verified
2026-10-11).
\item College Board, \emph{AP Physics C (E\&M) CED}, Unit 13 ---
verified 2026-10-11 at \url{https://apstudents.collegeboard.org}.
\item E.~M. Purcell and D.~J. Morin, \emph{Electricity and Magnetism},
chs.~7--8 --- induction and energy with the field bookkeeping.
\item PhET, \emph{Faraday's Law} interactive --- move the magnet, watch
the flux; at \url{https://phet.colorado.edu} (verified 2026-10-11).
\item Khan Academy, \emph{Faraday's law} ---
verified 2026-10-11 at \url{https://www.khanacademy.org}
(science $\to$ physics $\to$ E\&M).
\end{itemize}

\section*{Additional exercises}
\addcontentsline{toc}{section}{Additional exercises}

\begin{exercise}[easy]
A loop's flux drops from $0.4$ to $0.1$\,Wb in $0.2$\,s: find
$|\mathcal E|$. (AP-C)
\end{exercise}
\begin{exercise}[easy]
State the induced current's direction for a loop as a magnet's north
pole approaches; recede instead. (AP-C)
\end{exercise}
\begin{exercise}[easy]
A $200$\,mH inductor carries $3$\,A: find its stored energy. (AP-C)
\end{exercise}
\begin{exercise}[easy]
An RL circuit has $\tau=2$\,ms with $L=40$\,mH: find $R$. (F=ma)
\end{exercise}
\begin{exercise}[easy]
A $50$-turn coil of area $0.01$\,m$^{2}$ flips $180^{\circ}$ in a
$0.5$\,T field in $0.1$\,s: find the average EMF. (AP-C)
\end{exercise}
\begin{exercise}[medium]
Complete the sliding-rod audit: find the terminal rod speed for a given
pulling force. (AP-C FR-style)
\end{exercise}
\begin{exercise}[medium]
A loop enters a field region twice its width: sketch $I(t)$ over the
whole transit. (F=ma)
\end{exercise}
\begin{exercise}[medium]
Two coaxial loops: estimate the mutual inductance when the small loop
sits on the big one's axis. (AP-C)
\end{exercise}
\begin{exercise}[medium]
An RLC rings at $99\%$ of $\omega_0$: find $R/L$. (AP-C)
\end{exercise}
\begin{exercise}[medium]
Why does a transformer's secondary see alternating EMF but no net
charge transfer? Reason via mutual inductance. (engineering)
\end{exercise}
\begin{exercise}[medium]
An inductor and capacitor of given $L,C$: find the oscillation
frequency and the peak current for initial charge $Q_0$. (AP-C)
\end{exercise}
\begin{exercise}[hard]
A magnet falls through a conducting tube: argue the terminal velocity
exists and estimate it via energy balance. (USAPhO-style)
\end{exercise}
\begin{exercise}[hard]
A loop rotates at $\omega$ in a uniform field: derive
$\mathcal E(t)=NBA\omega\sin\omega t$ and explain the generator's
waveform. (APhO-style)
\end{exercise}
\begin{exercise}[medium]
Show $u_B+u_E$ inside a charging capacitor accounts for the energy
inflow. (concept)
\end{exercise}
\begin{exercise}[medium]
A betatron accelerates electrons: explain why the induced $\mathbf E$
direction stays fixed as $B$ ramps. (concept)
\end{exercise}
